\setcounter{chapCompteur}{5}
\ChapNumber{}{\hspace*{0.8cm}\LARGE \textbf{Sources et tests}}{\hspace*{0.8cm}\small Une source, deux façons de l'utiliser}
\section{Conclusion}

\begin{frame}[t]{Ce qu'il faut retenir}{Pour les développeurs : structure, tests, licence}
\vfill
{\relscale{0.9}
\begin{mybox3}[yellow!8!white][Framajaune][Framajaune]{\textbf{Points clés}}
\begin{itemize}\itemsep5pt
\item Source unique : \texttt{src/dim2dir/core.py}, commune au paquet et au fichier autonome.
\item \texttt{tools/sync\_standalone.py} synchronise la copie autonome \texttt{Dim2Dir.py}.
\item Tests : \texttt{pytest}, API, limites, exports et parité avec la version 1.
\item Licence : GNU AGPLv3, avec possibilité d'utiliser une version ultérieure.
\item Projet parent : \href{https://github.com/Philippe-AMUZUGA/plex2-doe-explorer}{dépôt PLEX2 DOE Explorer}~\cite{plex2}.
\end{itemize}
\end{mybox3}
Modifier la source unique, synchroniser la copie, puis exécuter les tests.}
\note{Le rendu par défaut reste protégé par les contrôles de parité. Tests : python -m pytest -q depuis la racine ; synchronisation : python tools/sync_standalone.py --check. Les tests Excel COM exigent Windows et Excel.}
\vfill
\end{frame}
